%============================================================
% PREÁMBULO OFICIAL FLUIDODINAMICA_BASE
%============================================================

\documentclass[12pt, a4paper]{article}

%------------------------------------------------------------
% Idioma, codificación y tipografía
%------------------------------------------------------------
\usepackage[spanish]{babel}
\usepackage[utf8]{inputenc}
\usepackage[T1]{fontenc}
\usepackage{lmodern}        % Tipografía moderna, clara y apta para PDF
\usepackage{microtype}      % Optimiza microtipografía (ideal para libros y apuntes)

%------------------------------------------------------------
% Matemática
%------------------------------------------------------------
\usepackage{amsmath, amssymb, amsfonts}
\usepackage{bm}             % Vectores y tensores en negrita: \bm{u}
\usepackage{mathtools}      % Extensiones de amsmath
\usepackage{siunitx}        % Notación correcta de unidades SI
\sisetup{
  locale = DE,            % usa coma decimal si querés; cambiar a US para punto
  per-mode = symbol,
  exponent-product = \cdot,
  separate-uncertainty = true
}

%------------------------------------------------------------
% Gráficos y figuras
%------------------------------------------------------------
\usepackage{graphicx}
\usepackage{float}
\usepackage{subcaption}     % Figuras con subfiguras
\usepackage{wrapfig}

%------------------------------------------------------------
% Diseño de página
%------------------------------------------------------------
\usepackage{geometry}
\geometry{
  a4paper,
  left=25mm,
  right=25mm,
  top=25mm,
  bottom=25mm
}

\usepackage{setspace}       % Control de espaciado (opcional)
\onehalfspacing             % Espaciado 1.5 (queda excelente para apuntes)

%------------------------------------------------------------
% Tablas profesionales
%------------------------------------------------------------
\usepackage{booktabs}       % Líneas horizontales profesionales
\usepackage{multirow}
\usepackage{tabularx}
\usepackage{longtable}

%------------------------------------------------------------
% Enlaces
%------------------------------------------------------------
\usepackage[hidelinks]{hyperref}  % Enlaces sin recuadros

%------------------------------------------------------------
% Colores y cajas
%------------------------------------------------------------
\usepackage{xcolor}
\usepackage{tcolorbox}
\tcbset{
  colback = white,
  colframe = black,
  arc = 2mm,
  boxrule = 0.6pt
}

%------------------------------------------------------------
% Ambientes personalizados
%------------------------------------------------------------

% Entorno para ejercicios
\newenvironment{ejercicio}[1][]{
  \begin{tcolorbox}[title={\textbf{Ejercicio #1}}, colback=gray!5]
}{
  \end{tcolorbox}
}

% Entorno para soluciones
\newenvironment{solucion}{
  \begin{tcolorbox}[title={\textbf{Solución}}, colback=blue!3]
}{
  \end{tcolorbox}
}

% Para dejar comentarios o notas al margen
\newcommand{\nota}[1]{\textcolor{blue}{\textbf{#1}}}

%============================================================
% FIN DEL PREÁMBULO
%============================================================
\begin{document}


%============================================================
% Ejercicio: Medidor Venturi en aire (flujo compresible moderado)
%============================================================

\subsection*{Ejercicio X: Medidor Venturi en una línea de aire}

En una cañería de aire de $2''$ Sch 40, que opera a $25\,^\circ\text{C}$, se instala 
un tubo Venturi como medidor de caudal. Un manómetro ubicado inmediatamente antes del 
Venturi indica una presión estática de $P_1 = 1{,}00\,\text{bar}$ y otro manómetro 
en la garganta indica $P_2 = 0{,}75\,\text{bar}$. 

El diámetro interno de la cañería de $2''$ Sch 40 es 
$D = 2{,}067'' \simeq 0{,}0525\,\text{m}$, mientras que el diámetro de la garganta es 
$d = 3{,}0\,\text{cm} = 0{,}030\,\text{m}$. Se dispone de la siguiente correlación para 
el coeficiente de descarga del Venturi:
%
\[
C_D = 0{,}986 - 0{,}196\,\beta^{4{,}5},
\qquad
\beta = \frac{d}{D}.
\]
%
Suponer que el aire se comporta como gas ideal y que la densidad puede 
aproximarse con un valor medio adecuado a las condiciones de operación. 
Utilizar, además, un factor de expansibilidad $Y$ representativo del efecto de 
compresibilidad.

\medskip

\noindent
\textbf{Datos adicionales:}
\begin{itemize}
  \item Aire: $\gamma = c_p/c_v = 1{,}4$, $\mu = 1{,}8\times 10^{-5}\,\text{Pa\,s}$,
        $T = 25\,^\circ\text{C} \simeq 298\,\text{K}$.
  \item Presiones: $P_1 = 1{,}00\,\text{bar}$, $P_2 = 0{,}75\,\text{bar}$.
        Por tanto, $\Delta P = P_1 - P_2 = 0{,}25\,\text{bar} = 2{,}5\times 10^4\,\text{Pa}$.
\end{itemize}

\noindent
\textbf{Se pide:} determinar el flujo másico de aire $\dot m$ en $\text{kg/h}$.

%------------------------------------------------------------
\subsubsection*{1. Datos, incógnitas y unidades}

\begin{itemize}
  \item Diámetro cañería (sección 1): $D = 0{,}0525\,\text{m}$.
  \item Diámetro garganta (sección 2): $d = 0{,}030\,\text{m}$.
        Relación de diámetros:
        \[
        \beta = \frac{d}{D} = \frac{0{,}030}{0{,}0525} \approx 0{,}571.
        \]
  \item Diferencia de presión entre tomas:
        \[
        \Delta P = P_1 - P_2 = 2{,}5\times 10^4\,\text{Pa}.
        \]
  \item Temperatura: $T \simeq 298\,\text{K}$.
  \item Propiedades del aire (valor medio tabulado): 
        \[
        \rho \simeq 2{,}3\,\text{kg/m}^3
        \]
        evaluada a una presión media $P_\text{med} \approx 0{,}875\,\text{bar}$ 
        y $T \simeq 298\,\text{K}$.
  \item Se adopta un factor de expansibilidad típico $Y \simeq 0{,}92$.
  \item Incógnita: flujo másico de aire $\dot m$ en $\text{kg/s}$ y su valor en $\text{kg/h}$.
\end{itemize}

%------------------------------------------------------------
\subsubsection*{2. Hipótesis físicas y termodinámicas}

\begin{enumerate}
  \item Flujo estacionario, monofásico, de aire.
  \item El aire se modela como gas ideal con propiedades constantes en el tramo considerado.
  \item El Venturi se trata como un medidor por restricción de área, de acuerdo al teórico:
        se parte del modelo ideal (continuidad + Bernoulli) y se introducen 
        un coeficiente de descarga $C_D(\beta)$ y un factor de expansibilidad $Y$ para gas.
  \item Se desprecia la diferencia de cota entre secciones 
        ($z_1 \simeq z_2$) frente a los efectos de velocidad y presión.
  \item La densidad se aproxima como constante e igual a un valor medio sujeto a verificación
        posterior (a través del número de Mach).
\end{enumerate}

%------------------------------------------------------------
\subsubsection*{3. Interpretación del sistema}

El Venturi se inserta en una línea de aire en régimen permanente. La sección 1 corresponde 
a la entrada del Venturi (diámetro $D$, presión estática $P_1$, velocidad $v_1$), mientras
que la sección 2 corresponde a la garganta (diámetro $d$, presión estática $P_2$, 
velocidad $v_2$). La reducción de área causa un aumento de la velocidad de escoamiento 
en la garganta y una correspondiente caída de presión, medida por la diferencia 
$\Delta P = P_1 - P_2$. Dicha diferencia de presión se utiliza para inferir el caudal 
mediante la ecuación reducida de Venturi.

%------------------------------------------------------------
\subsubsection*{4. Ecuaciones gobernantes}

\paragraph{Continuidad.}
Para flujo estacionario con densidad aproximadamente uniforme entre 1 y 2, se tiene
\[
Q = v_1 A_1 = v_2 A_2,
\]
donde $A_1$ y $A_2$ son las áreas de las secciones 1 y 2, respectivamente. A partir de 
esta relación se puede escribir
\[
v_1 = \frac{Q}{A_1},\qquad v_2 = \frac{Q}{A_2}, 
\]
y, usando $\beta^2 = A_2/A_1$, se cumple
\[
v_1 = \beta^2\,v_2.
\]

\paragraph{Ecuación de Bernoulli ideal.}
Entre las secciones 1 y 2, despreciando diferencias de cota y pérdidas distribuidas entre 
las tomas de presión, el Bernoulli ideal establece
\[
P_1 + \tfrac{1}{2}\rho v_1^2 = P_2 + \tfrac{1}{2}\rho v_2^2.
\]
Sustrayendo $P_2 + \tfrac{1}{2}\rho v_1^2$ de ambos lados y usando $v_1 = \beta^2 v_2$:
\[
P_1 - P_2 = \tfrac{1}{2}\rho\left(v_2^2 - v_1^2\right)
          = \tfrac{1}{2}\rho v_2^2 \left(1 - \beta^4\right).
\]
Por tanto, la velocidad ideal en la garganta viene dada por
\[
v_{2,\text{ideal}} = \sqrt{\frac{2\,\Delta P}{\rho(1-\beta^4)}}.
\]

\paragraph{Modelo real del Venturi para gas.}
Siguiendo el teórico de clase (medidores por restricción de área), las desviaciones del 
modelo ideal se recogen mediante:
\begin{itemize}
  \item un coeficiente de descarga $C_D(\beta)$,
  \item un factor de expansibilidad $Y$ para tener en cuenta la compresibilidad del aire.
\end{itemize}
La ecuación empírica para el flujo másico a través del Venturi se escribe entonces como
\[
\boxed{
\dot m = C_D\,Y\,A_2\,\sqrt{\frac{2\,\rho\,\Delta P}{1-\beta^4}}
}
\]
donde $\dot m$ es el caudal másico real.

%------------------------------------------------------------
\subsubsection*{5. Cierre del modelo}

\begin{itemize}
  \item Se utiliza la correlación dada para el coeficiente de descarga:
        \[
        C_D = 0{,}986 - 0{,}196\,\beta^{4{,}5}.
        \]
  \item Se toma un valor medio de densidad del aire 
        $\rho \simeq 2{,}3\,\text{kg/m}^3$, evaluado a una presión media 
        $P_\text{med} \approx 0{,}875\,\text{bar}$ y $T=298\,\text{K}$.
  \item Se adopta un factor de expansibilidad representativo 
        $Y \simeq 0{,}92$ para la caída de presión indicada.
  \item Con estos valores, junto con $D$, $d$ y $\Delta P$, el modelo queda completamente 
        determinado para el cálculo de $\dot m$.
\end{itemize}

%------------------------------------------------------------
\subsubsection*{6. Desarrollo simbólico}

La expresión simbólica final se ha obtenido en el apartado anterior:
\[
\dot m(\Delta P) 
= C_D(\beta)\,Y\,A_2\,
\sqrt{\frac{2\,\rho\,\Delta P}{1-\beta^4}}.
\]
No se requieren transformaciones adicionales; se procede a la sustitución numérica.

%------------------------------------------------------------
\subsubsection*{7. Sustitución numérica ordenada}

\paragraph{a) Geometría.}

Relación de diámetros:
\[
\beta = \frac{d}{D} 
      = \frac{0{,}030}{0{,}0525}
      = 0{,}571.
\]

Cálculo de $\beta^{4{,}5}$:
\[
\beta^{4{,}5} \approx (0{,}571)^{4{,}5} \approx 0{,}080.
\]

Coeficiente de descarga:
\[
C_D = 0{,}986 - 0{,}196\,\beta^{4{,}5}
    \approx 0{,}986 - 0{,}196\times 0{,}080
    \approx 0{,}97.
\]

Área de la garganta:
\[
A_2 = \frac{\pi d^2}{4}
    = \frac{\pi(0{,}030)^2}{4}
    = 7{,}068\times 10^{-4}\,\text{m}^2.
\]

Término geométrico en el denominador:
\[
1-\beta^4 = 1 - (0{,}571)^4
          \approx 1 - 0{,}106
          = 0{,}894.
\]

\paragraph{b) Término bajo la raíz.}

Se evalúa
\[
\frac{2\,\rho\,\Delta P}{1-\beta^4}
= \frac{2\cdot 2{,}3\cdot 2{,}5\times 10^4}{0{,}894}
\approx 1{,}29\times 10^5\ \text{(m}^2/\text{s}^2).
\]

La raíz cuadrada resulta
\[
\sqrt{\frac{2\,\rho\,\Delta P}{1-\beta^4}}
\approx \sqrt{1{,}29\times 10^5}
\approx 3{,}59\times 10^2\,\text{m/s}.
\]

\paragraph{c) Cálculo del flujo másico.}

Sustituyendo en la expresión reducida:
\[
\dot m 
= C_D\,Y\,A_2\,
  \sqrt{\frac{2\,\rho\,\Delta P}{1-\beta^4}}.
\]

Con $C_D \approx 0{,}97$, $Y \approx 0{,}92$, 
$A_2 = 7{,}068\times 10^{-4}\,\text{m}^2$ y la raíz calculada:
\[
\dot m 
\approx 0{,}97\cdot 0{,}92\cdot 7{,}068\times 10^{-4}
       \cdot 3{,}59\times 10^2.
\]

Primero:
\[
0{,}97\cdot 0{,}92 \approx 0{,}892,
\quad
0{,}892\cdot 7{,}068\times 10^{-4}
\approx 6{,}30\times 10^{-4}.
\]

Luego:
\[
\dot m 
\approx 6{,}30\times 10^{-4}\cdot 3{,}59\times 10^2
\approx 2{,}27\times 10^{-1}\,\text{kg/s}.
\]

Por tanto,
\[
\boxed{\dot m \approx 0{,}227\,\text{kg/s}}.
\]

Conversión a $\text{kg/h}$:
\[
\dot m_\text{(kg/h)} 
= 0{,}227\,\frac{\text{kg}}{\text{s}}\cdot 3600\,\frac{\text{s}}{\text{h}}
\approx 8{,}16\times 10^2\,\text{kg/h}.
\]

Es decir,
\[
\boxed{\dot m \approx 8{,}2\times 10^2\,\text{kg/h}}.
\]

%------------------------------------------------------------
\subsubsection*{8. Verificación física y adimensional}

\paragraph{Velocidad y número de Reynolds.}

La velocidad en la garganta se obtiene como
\[
v_2 = \frac{\dot m}{\rho A_2}
    \approx \frac{0{,}227}{2{,}3\cdot 7{,}068\times 10^{-4}}
    \approx 1{,}4\times 10^2\,\text{m/s}.
\]

El número de Reynolds asociado a la garganta es
\[
\text{Re}_2 
= \frac{\rho v_2 d}{\mu}
\approx \frac{2{,}3\cdot 1{,}4\times 10^2\cdot 0{,}030}{1{,}8\times 10^{-5}}
\sim 5\times 10^5,
\]
lo que indica un régimen fuertemente turbulento, coherente con el uso de un 
coeficiente de descarga prácticamente constante.

\paragraph{Número de Mach.}

La velocidad del sonido en aire a $25\,^\circ\text{C}$ es aproximadamente 
$a \simeq 346\,\text{m/s}$, por lo que
\[
\text{Ma}_2 = \frac{v_2}{a}
            \approx \frac{140}{346}
            \approx 0{,}4.
\]

Se trata de un número de Mach moderado: la compresibilidad no es despreciable, lo que 
justifica la introducción del factor de expansibilidad $Y<1$ en el modelo de medición.

%------------------------------------------------------------
\subsubsection*{9. Comentario final}

El modelo empleado sigue el esquema teórico de medidores por restricción de área: 
se parte de las ecuaciones de continuidad y Bernoulli ideal, se introduce la relación 
geométrica $1-\beta^4$ y luego se corrige el caudal ideal mediante el coeficiente de 
descarga $C_D(\beta)$ y el factor de expansibilidad $Y$ para flujo compresible moderado
de aire. 

El valor obtenido, $\dot m \approx 0{,}23\,\text{kg/s}$ (del orden de $8\times 10^2\,\text{kg/h}$),
es razonable para un Venturi en una cañería de $2''$ con las diferencias de presión dadas 
y el régimen resultante (turbulento, $\text{Re}\sim 10^5$, $\text{Ma}\sim 0{,}4$) es coherente 
con el uso de las correlaciones experimentales de $C_D$ y $Y$ adoptadas en la asignatura.

\end{document}